\begin{frame}{확인 문제 (3.2)}
  \begin{enumerate}
    \item $\gamma = 0.9$, $\epsilon = 0.05$ 일 때 유효 시계 $H$ 를
      $H < \log(0.05)/\log(0.9)$ 로 직접 계산. 로그를 쓰지 않고
      $0.9^{28}$, $0.9^{29}$ 를 직접 곱해서도 검증. 이 $H = 28$ 이
      ``에이전트는 28스텝 이후의 보상을 5\% 이하로
      취급한다''는 뜻과 일치하는지 확인.
    \item 매 스텝 보상 $R = 1$ 인 환경에서, $\gamma = 0.9$ 와
      $\gamma = 0.99$ 일 때 각각 ``한 에피소드에서 받을 수
      있는 \textbf{최대} 리턴''($1/(1-\gamma)$)을 계산. 이 차이가
      ``$\gamma = 0.99$ 를 쓰면 $\gamma = 0.9$ 때보다
      \textbf{상대적으로} 더 먼 미래의 보상이 중요하다''는 뜻인지
      한 문장으로 설명.
    \item (개념) ``할인율 0.95는 ``5\% 이자율''''이라는 비유.
      금융에서 ``1년 뒤 1원''의 현재가치 $1/(1+r)$ 과 비교해,
      강화학습의 $\gamma$ 가 ``이자율''과 \textbf{같은 방향}인지,
      \textbf{다르다}는 점(예: $\gamma = 1$ 이 ``이자율 0''이
      아니라 ``미래 보상을 무한히 중시''하는 극단)을 각 한 문장으로.
  \end{enumerate}
\end{frame}
